% Part: lambda-calculus
% Chapter: introduction
% Section: overview

\documentclass[../../../include/open-logic-section]{subfiles}

\begin{document}

\olfileid[pt]{lam}{int}{ovr}
\olsection{Visão geral}

O cálculo lambda foi originalmente concebido por Alonzo Church, no início
da década de 1930, como base para a lógica construtiva, e \emph{não} como
modelo das funções computáveis. Logo, porém, mostrou-se que ele era
equivalente a outras definições de computabilidade, como as funções
computáveis por Turing e as funções recursivas parciais. O fato de isso ter
causado inicialmente uma pequena surpresa torna a caracterização ainda mais
interessante.

A notação lambda é uma maneira conveniente de se referir diretamente a uma
função por meio de uma expressão simbólica que a define, em vez de lhe dar
um nome. Em vez de dizer ``seja $f$ a função definida por
$f(x) = x + 3$'', podemos dizer ``seja $f$ a função
$\lambd[x][(x + 3)]$''. Em outras palavras, $\lambd[x][(x+3)]$ é apenas um
\emph{nome} para a função que soma três a seu argumento. Nessa expressão,
$x$ é uma variável muda, ou um marcador de posição: a mesma função pode
igualmente ser denotada por $\lambd[y][(y + 3)]$. A notação também funciona
na presença de outros parâmetros. Por exemplo, suponha que $g(x, y)$ seja
uma função de duas variáveis e $k$ um número natural. Então
$\lambd[x][g(x,k)]$ é a função que leva cada~$x$ em~$g(x, k)$.

Essa maneira de definir uma função a partir de uma expressão simbólica é
conhecida como \emph{abstração lambda}. O outro lado da abstração lambda é a
\emph{aplicação}: supondo que tenhamos uma função $f$ (digamos, definida nos
números naturais), podemos \emph{aplicá-la} a qualquer valor, como 2. Na
notação convencional, é claro, escrevemos~$f(2)$ para o resultado.

O que acontece quando combinamos abstração lambda com aplicação? A
expressão resultante pode ser simplificada ``inserindo'' o aplicando no
lugar da variável abstraída. Por exemplo,
\[
(\lambd[x][(x + 3)])(2)
\]
pode ser simplificada para~$2 + 3$.

Até aqui, nada fizemos além de introduzir novas notações para noções
convencionais. O cálculo lambda, porém, representa um afastamento mais
radical do ponto de vista da teoria dos conjuntos. Nesse quadro:
\begin{enumerate}
\item Tudo denota uma função.
\item Funções podem ser definidas por abstração lambda.
\item Qualquer coisa pode ser aplicada a qualquer outra coisa.
\end{enumerate}
Por exemplo, se $F$ é um termo no cálculo lambda, sempre se supõe que $F(F)$
tenha sentido. Esse quadro liberal é conhecido como cálculo lambda
\emph{não tipado}, onde ``não tipado'' significa ``sem restrição sobre o que
pode ser aplicado a quê''.

\begin{digress}
Há também um cálculo lambda \emph{tipado}, que é uma variação importante da
versão não tipada. Embora o cálculo lambda tipado seja, em muitos aspectos,
semelhante ao não tipado, ele é muito mais fácil de conciliar com um quadro
clássico da teoria dos conjuntos e tem algumas propriedades bastante
diferentes.

A pesquisa sobre o cálculo lambda mostrou-se central para a ciência da
computação teórica e para o projeto de linguagens de programação. LISP,
projetada por John McCarthy na década de 1950, é um dos primeiros exemplos
de linguagem influenciada por essas ideias.
\end{digress}

\end{document}

